% !TeX root = ../main.tex

% ==================================================
% 2.9. super()와 MRO — Python은 다중 상속에서 메서드를 어떤 순서로 탐색하는가?
% ==================================================

\subsection{\zsh{super()}와 MRO - Python은 다중 상속에서 메서드를 어떤 순서로 탐색하는가?}

2.8에서는 상속이 부모 클래스의 메서드를 자식 클래스에 복사하는 것이 아니라, 클래스 사이에 속성 탐색 관계를 형성하는 것이라고 학습했다.

이번에는 그 관계를 조금 더 엄밀하게 살펴볼 것이다.

특히 다음 코드를 이해하는 것이 목표다.

\begin{dhkcode}
class Dog(Animal):
    def __init__(self, name):
        super().__init__(name)
\end{dhkcode}

우리는 흔히 \zsh{super().__init__()}을 부모 클래스의 \zsh{__init__()}을 호출하는 문법이라고 설명한다.

단일 상속에서는 대체로 그렇게 이해해도 결과가 맞다.

하지만 Python의 실제 동작 원리는 조금 다르다.

\zsh{super()}는 부모 클래스를 직접 지정하는 것이 아니라, MRO를 기준으로 현재 클래스 다음에 위치한 클래스부터 속성을 탐색하는 메커니즘이다.

이 차이는 다중 상속에서 매우 중요해진다.

\subsubsection{MRO란 무엇인가?}

MRO는 \zsh{Method Resolution Order}의 약자다.

Python이 클래스의 속성이나 메서드를 탐색할 때 사용하는 클래스들의 선형적인 탐색 순서를 의미한다.

다음 코드를 보자.

\begin{dhkcode}
class Animal:
    def speak(self):
        print("Animal")


class Dog(Animal):
    pass
\end{dhkcode}

\zsh{Dog}의 MRO를 확인해 보자.

\begin{dhkcode}
print(Dog.__mro__)
\end{dhkcode}

실행 결과:

\begin{dhkzsh}
(<class '__main__.Dog'>,
 <class '__main__.Animal'>,
 <class 'object'>)
\end{dhkzsh}

이 결과는 Python이 \zsh{Dog} 인스턴스에서 클래스 속성을 탐색할 때 다음 순서를 사용한다는 의미다.

\begin{dhkgraph}[uniform]
node n1 (1,2) [accent] :: Dog // 현재 클래스
node n2 (2,2) :: Animal // 부모 클래스
node n3 (3,2) :: object // 최상위 기본 클래스
node s1 (1,1) :: 1.
node s2 (2,1) :: 2.
node s3 (3,1) :: 3.
edge n1 -> n2
edge n2 -> n3

caption :: Dog.__mro__에 저장된 클래스 탐색 순서
\end{dhkgraph}

MRO는 다음처럼도 확인할 수 있다.

\begin{dhkcode}
pritn(Dog.mro())
\end{dhkcode}

두 표현식의 차이는 반환 형식이다.

\begin{dhktable}
|표현식|반환값|
|:---|:---|
|Dog.__mro__|클래스들의 튜플|
|Dog.mro()|클래스들의 리스트|
\end{dhktable}

일반적인 클래스에서는 두 표현식으로 동일한 탐색 순서를 확인할 수 있다.

\begin{dhknote}
\begin{dhksmry}[centered]
중요한 구분: MRO는 클래스 탐색 순서다
\end{dhksmry}
\end{dhknote}

2.7에서 배운 인스턴스 속성 탐색에는 디스크립터와 인스턴스 \zsh{__dict__}의 우선순위가 있었다.

MRO는 그 규칙을 대체하는 것이 아니다.

MRO는 클래스 계층에서 속성을 탐색해야 할 때 어떤 클래스를 먼저 살펴볼지 결정한다.

따라서 다음 두 개념을 구분해야 한다.

\begin{dhksmry}
- 속성 탐색 규칙: 데이터 디스크립터, 인스턴스 딕셔너리 등의 우선순위

- MRO : 클래스와 상속 계층의 탐색 순서
\end{dhksmry}

이 둘이 결합되어 실제 속성 접근 결과가 결정되는 것이다.

\subsubsection{\zsh{super()}는 무엇을 반환하는가?}

다음 코드를 보자.

\begin{dhkcode}
class Animal:
    def speak(self):
        print("Animal")


class Dog(Animal):
    def speak(self):
        super().speak()
        print("Dog")

dog = Dog()
dog.speak()
\end{dhkcode}

실행 결과:

\begin{dhkzsh}
Animal
Dog
\end{dhkzsh}

이때 다음 표현식에 주목한다.

\begin{dhkcode}
super()
\end{dhkcode}

이 표현식은 \zsh{Animal} 클래스 객체를 반환하는 것이 아니다.

MRO를 기준으로 속성을 탐색하도록 구성된 \zsh{super} 프록시 객체를 반환한다.

이를 직접 확인해 보자.

\begin{dhkcode}
class Dog(Animal):
    def inspect(self):
        proxy = super()

        print(type(proxy))
        print(proxy)

dog = Dog()
dog.inspect()
\end{dhkcode}

출력 예시:

\begin{dhkzsh}
<class 'super'>
<super: <class 'Dog'>, <Dog object>>
\end{dhkzsh}

여기서 \zsh{super()} 객체는 두 가지 중요한 정보를 이용한다.

\begin{dhksmry}
1. 탐색을 시작할 기준 클래스

2. 속성 탐색에 사용할 인스턴스 또는 클래스
\end{dhksmry}

이 예제에서는 기준 클래스가 \zsh{Dog}, 연결된 객체가 \zsh{dog}이다.

따라서 \zsh{super().speak()}은 \zsh{Dog}의 MRO에서 \zsh{Dog}다음 위치부터 \zsh{speak}을 찾는 것이다.

\begin{dhkdiagram}[in=6]
node n1 (1,1) :: Dog.__mro__
node n2 (3,1) :: Dog
node n3 (3,3) :: Animal
node n4 (3,5) :: object
node n5 (5,3) :: super().speak()

edge n2 -> n3
edge n3 -> n4
edge n5 -> n3
\end{dhkdiagram}

단일 상속에서는 그 다음 클래스가 보통 직접적인 부모 클래스이기 때문에 \zsh{super()}가 부모 메서드를 호출하는 것처럼 보이는 것이다.

\subsubsection{\zsh{super()}의 두 인수 형태}

Python에서는 다음 표현식을 사용할 수 있다.

\begin{dhkcode}
super(Dog, dog)
\end{dhkcode}

일반적인 인스턴스 메서드 안에서 사용하는

\begin{dhkcode}
super()
\end{dhkcode}

는 적절한 문맥에서 다음과 같은 효과를 가진다.

\begin{dhkcode}
super(Dog, self)
\end{dhkcode}

여기서 첫 번째 인수는 MRO에서 탐색을 건너뛸 기준 클래스다.

두 번째 인수는 그 탐색에 사용할 인스턴스다.

직접 확인해 보자.

\begin{dhkcode}
class Animal:
    def speak(self):
        print("Animal")


class Dog(Animal):
    def speak(self):
        super(Dog, self).speak()
        print("Dog")

dog = Dog()
dog.speak()
\end{dhkcode}

실행 결과:

\begin{dhkzsh}
Animal
Dog
\end{dhkzsh}

두 인수를 명시적으로 전달해도 같은 결과가 나온다.

하지만 일반적인 인스턴스 메서드에서는 다음 표현식을 권장한다.

\begin{dhkcode}
super().speak()
\end{dhkcode}

클래스 이름을 직접 작성하지 않아도 되므로 클래스 이름 변경이나 상속 구조 변경에 더 유연하기 때문이다.

\begin{dhknote}
\begin{dhksmry}[centered]
super()가 무조건 상위 클래스를 의미할까?
\end{dhksmry}
\end{dhknote}

아니다.

정확하게는 지정한 기준 클래스 다음부터 MRO를 따라 탐색하는 객체라고 이해해야 한다.

이 사실은 다중 상속에서 분명해진다.

\subsubsection{다중 상속 - 부모 클래스가 여러 개인 경우}

Python은 다중 상속(Multiple Inheritance)을 지원한다.

즉, 하나의 클래스가 여러 클래스를 동시에 상속할 수 있다.

다음 코드를 보자.

\begin{dhkcode}
class A:
    def speak(self):
        print("A")


class B:
    def speak(self):
        print("B")


class C(A, B):
    pass
\end{dhkcode}

\zsh{C}는 \zsh{A}와 \zsh{B}를 모두 상속한다.

\begin{dhkcode}
print(C.__bases__)
\end{dhkcode}

실행 결과:

\begin{dhkzsh}
(<class '__main__.A'>, <class '__main__.B'>)
\end{dhkzsh}

그렇다면 \zsh{C}의 인스턴스에서 \zsh{speak()}을 호출하면 어떤 메서드가 실행될까?

\begin{dhkcode}
c = C()
c.speak()
\end{dhkcode}

실행 결과:

\begin{dhkzsh}
A
\end{dhkzsh}

이유는 MRO를 보면 알 수 있다.

\begin{dhkcode}
print([cls.__name__ for cls in C.__mro__])
\end{dhkcode}

실행 결과:

\begin{dhkzsh}
['C', 'A', 'B', 'object']
\end{dhkzsh}

즉 클래스 탐색 순서는 다음과 같다.

\begin{dhkgraph}
node n1 (1,2) [accent] :: C
node n2 (2,1) :: A
node n3 (2,3) :: B
node n4 (3,2) :: MRO: C->A->B->object

edge n1 -- n2
edge n1 -- n3

caption :: 상속 구조는 분기하지만, MRO는 하나의 선형 순서를 제공한다.
\end{dhkgraph}

이 예제에서는 \zsh{A}와 \zsh{B}가 모두 \zsh{speak()}을 정의했다.

Python은 MRO에서 먼저 등장하는 \zsh{A}의 메서드를 선택한다.

여기서 중요한 점은 다음이다.

MRO는 단순히 부모 클래스의 목록이 아니라, 상속 계층 전체를 하나의 일관된 탐색 순서로 정렬한 결과다.

\subsubsection{다이아몬드 상속 - 왜 MRO가 필요한가?}

이제 조금 더 복잡한 구조를 살펴본다.

\begin{dhkcode}
class A:
    def method(self):
        print("A")


class B(A):
    pass


class C(A):
    pass


class D(B, C):
    pass
\end{dhkcode}

이 구조를 그림으로 나타내면 다음과 같다.

\begin{dhkgraph}
node n1 (1,2) :: A
node n2 (2,1) :: B
node n3 (2,3) :: C
node n4 (3,2) :: D
node n5 (4,2) :: MRO: D -> B -> C -> A -> object

edge n1 -- n2
edge n1 -- n3
edge n2 -- n4
edge n3 -- n4

caption :: 다이아몬드 상속: B와 C가 같은 조상 A를 상속하고, D가 두 클래스를 상속한다.
\end{dhkgraph}

이 구조를 다이아몬드 상속(Diamond Inheritance)이라고 한다.

여기서 문제가 생길 수 있다.

\zsh{D}는 \zsh{B}를 통해 \zsh{A}에 도달할 수도 있고, \zsh{C}를 통해 \zsh{A}에 도달할 수도 있다.

그렇다면 \zsh{A}는 두 번 탐색해야 할까?

Python은 MRO를 사용해 이 문제를 해결한다.

\begin{dhkcode}
print([cls.__name__ for cls in D.__mro__])
\end{dhkcode}

실행 결과:

\begin{dhkzsh}
['D', 'B', 'C', 'A', 'object']
\end{dhkzsh}

여기서 \zsh{A}는 한 번만 등장한다.

이처럼 MRO는 다중 상속에서 클래스 탐색 순서를 일관되게 유지하는 역할을 한다.

\subsubsection{C3 선형화 - MRO는 어떤 규칙으로 결정되는가?}

Python 3에서는 일반적인 클래스의 MRO를 결정할 때 C3 선형화(C3 Linearization) 알고리즘을 사용한다.

C3는 단순히 깊이 우선 탐색(DFS)이나 너비 우선 탐색(BFS)을 수행하는 알고리즘이 아니다.

다중 상속에서 여러 제약을 만족하도록 클래스들을 선형 순서로 배치하는 알고리즘이다.

핵심 원칙은 다음과 같다.

\begin{dhksmry}
1. 자식 클래스는 부모 클래스보다 먼저 탐색된다.

2. 클래스 선언에서 지정한 직접 부모 클래스들의 상대적 순서를 존중한다.

3. 부모 클래스들이 이미 가진 MRO의 상대적 순서를 유지한다.

4. 이 조건들을 만족하는 일관된 순서가 존재하지 않으면 클래스 생성을 거부한다.
\end{dhksmry}

이를 수식으로 나타내면 다음과 같다.

\[
L[C] =
[C] +
\operatorname{merge}
\left(
L[B_1],L[B_2],\ldots,L[B_n],
[B_1,B_2,\ldots,B_n]
\right)
\]

여기서

\(L[C]\)는 클래스 \(C\)의 MRO

\(B_1,\ldots,B_n\)은 직접 부모 클래스

\(\operatorname{merge}\)는 C3의 병합 연산

이다.

\begin{dhknote}
\begin{dhksmry}[centered]
C3 병합 규칙
\end{dhksmry}
\end{dhknote}

병합 과정에서는 각 리스트의 첫 번째 원소를 후보로 살펴본다.

그 후보가 다른 리스트의 첫 번째가 아닌 나머지 부분에 등장하지 않으면 결과에 추가할 수 있다.

후보를 선택했다면 각 리스트에서 해당 원소를 제거하고 같은 과정을 반복한다.

이 규칙이 중요한 이유는 기존 부모 클래스들의 상대적 탐색 순서를 함부로 뒤집지 않기 때문이다.

\begin{dhknote}
\begin{dhksmry}[centered]
다이아몬드 구조에 적용하기
\end{dhksmry}
\end{dhknote}

앞의 예제를 다시 보자.

\begin{dhkcode}
class A:
    pass

class B(A):
    pass

class C(A):
    pass

class D(B, C):
    pass
\end{dhkcode}

각 클래스의 MRO는 다음과 같다.

\begin{dhksmry}
    A: A → object

    B: B → A → object

    C: C → A → object

    D: D → B → C → A → object
\end{dhksmry}

\zsh{D}의 MRO를 구성할 때는 다음 리스트들을 병합해야 한다.

\[
\operatorname{merge}
\left(
[B,A,object],
[C,A,object],
[B,C]
\right)
\]

첫 번째로 \zsh{B}를 선택할 수 있다.

\zsh{B}는 다른 리스트들의 나머지 부분에 등장하지 않기 때문이다.

그다음 \zsh{C}를 선택하고, 이어서 \zsh{A}, 마지막으로 \zsh{object}를 선택한다.

따라서 최종 결과는 다음과 같다.

\[
\boxed{
L[D] = [D,B,C,A,object]
}
\]

여기서 중요한 사실은 C3가 클래스 선언의 순서만 단순히 이어 붙이는 것이 아니라, 상속 계층 전체의 순서 제약을 만족하도록 병합한다는 것이다.

\begin{dhknote}
\begin{dhksmry}[centered]
일관된 순서를 만들 수 없는 경우
\end{dhksmry}
\end{dhknote}

\begin{dhkcode}
class A:
    pass

class B:
    pass

class X(A, B):
    pass

class Y(B, A):
    pass

class Z(X, Y):
    pass
\end{dhkcode}

이 코드는 \zsh{Z}를 정의할 때 \zsh{TypeError}가 발생한다.

왜냐하면 \zsh{X}는 \zsh{A}가 \zsh{B}보다 먼저 와야 한다고 요구하고,

\zsh{Y}는 \zsh{B}가 \zsh{A}보다 먼저 와야 한다고 요구하기 때문이다.

두 제약을 동시에 만족하는 MRO는 존재하지 않는다.

따라서 Python은 임의로 하나를 선택하는 대신 클래스 정의 자체를 거부한다.

이것이 C3 선형화의 일관성 조건이다.

\subsubsection{\zsh{super()}는 다중 상속에서 어떻게 동작하는가?}

이제 이번 학습에서 가장 중요한 예제로 들어가 보자.

\begin{dhkcode}
class A:
    def process(self):
        print("A")


class B(A):
    def process(self):
        print("B 시작")
        super().process()
        print("B 종료")


class C(A):
    def process(self):
        print("C 시작")
        super().process()
        print("C 종료")


class D(B, C):
    def process(self):
        print("D 시작")
        super().process()
        print("D 종료")
\end{dhkcode}

이제 다음 코드를 실행해 보자.

\begin{dhkcode}
d = D()
d.process()
\end{dhkcode}

실행 결과:

\begin{dhkzsh}
D 시작
B 시작
C 시작
A
C 종료
B 종료
D 종료
\end{dhkzsh}

왜 이런 순서가 나올까?

먼저 MRO를 확인하자.

\begin{dhkcode}
print([cls.__name__ for cls in D.__mro__])
\end{dhkcode}

실행 결과:

\begin{dhkzsh}
['D', 'B', 'C', 'A', 'object']
\end{dhkzsh}

여기서 중요한 부분은 \zsh{B.process()} 안에 있는 다음 코드다.

\begin{dhkcode}
super().process()
\end{dhkcode}

일반적으로 \zsh{B}의 부모는 \zsh{A}이므로 \zsh{A.process()}를 호출할 것이라고 생각하기 쉽다.

하지만 실제로는 \zsh{C.process()}가 호출된다.

왜 그럴까?

\zsh{super()}가 \zsh{B}의 직접 부모 클래스를 호출하는 것이 아니라, 실제 \zsh{d} 인스턴스의 MRO에서 \zsh{B} 다음에 있는 클래스를 탐색하기 때문이다.

\begin{dhkgraph}[uniform]
node n1 (1,1) :: D.process() // D 시작
node n2 (2,1) :: B.process() // B 시작
node n3 (3,1) [accent] :: C.process() // C 시작
node n4 (4,1) :: A.process() // A 출력

edge n1 -> n2
edge n2 -> n3
edge n3 -> n4

caption :: 각 super().process()는 동일한 인스턴스 d의 MRO에서 다음 구현을 찾아 호출한다.
\end{dhkgraph}

호출 순서를 조금 더 엄밀하게 표현하면 다음과 같다.

\begin{dhktable}
|실행 위치|	super()의 기준 클래스|	다음에 발견되는 메서드|
|:---|:---|:---|
|D.process()|	D|	B.process()|
|B.process()|	B|	C.process()|
|C.process()|	C|	A.process()|
\end{dhktable}

이 표에서 특히 두 번째 행이 중요하다.

\zsh{B}의 직접 부모는 \zsh{A}이지만, \zsh{D} 인스턴스의 MRO에서는 \zsh{B} 다음에 \zsh{C}가 위치한다.

따라서 \zsh{B.process()} 안의 \zsh{super()}가 \zsh{C.process()}를 호출하는 것이다.

\begin{dhknote}
\begin{dhksmry}[centered]
그렇다면 B() 인스턴스에서는?
\end{dhksmry}
\end{dhknote}

\begin{dhkcode}
b = B()
b.process()
\end{dhkcode}

실행 결과:

\begin{dhkzsh}
B 시작
A
B 종료
\end{dhkzsh}

왜 \zsh{C}를 거치지 않을까?

\zsh{b}의 타입은 \zsh{B}이고, MRO는 다음과 같기 때문이다.

\begin{dhkzsh}
B → A → object
\end{dhkzsh}

즉, 같은 \zsh{B.process()} 함수 객체라도 어떤 인스턴스의 MRO를 기준으로 실행되는지에 따라 \zsh{super()}의 다음 탐색 대상이 달라질 수 있다.

이것이 \zsh{super()}를 단순히 부모 메서드 호출이라고만 이해하면 안 되는 가장 중요한 이유다.

\subsubsection{협력적 다중 상속 - Cooperative Multiple inheritance}

앞의 예제에서는 각 클래스가 다음 코드를 사용했다.

\begin{dhkcode}
super().process()
\end{dhkcode}

이처럼 각 클래스가 MRO의 다음 구현으로 제어를 전달하도록 설계하는 방식을

\zsh{협력적 다중 상속(Cooperative Multiple Inheritance)}이라고 한다.

여기서 중요한 것은 메서드들이 서로 협력하도록 설계되어야 한다는 점이다.

예를 들어 어떤 클래스가 다음처럼 작성되어 있다면

\begin{dhkcode}
class B(A):
    def process(self):
        print("B")
        A.process(self)
\end{dhkcode}

\zsh{B}는 다음 MRO 항목을 찾지 않고 \zsh{A.process()}를 직접 호출한다.

이렇게 하면 다중 상속 구조에서 중간에 위치한 다른 클래스의 메서드를 건너뛸 수 있다.

반면 다음 코드는

\begin{dhkcode}
class B(A):
    def process(self):
        print("B")
        super().process()
\end{dhkcode}

실제 인스턴스의 MRO에 따라 다음 구현을 호출할 수 있다.

\begin{dhknote}
\begin{dhksmry}[centered]
하지만 super()만 사용하면 항상 안전할까?
\end{dhksmry}
\end{dhknote}

그렇지는 않다.

협력적 다중 상속에서는 다음 조건들도 중요하다.

\begin{dhksmry}
- 각 메서드가 다음 구현을 호출하도록 설계되어야 한다.

- 메서드들의 인수 전달 규약이 서로 호환되어야 한다.

- 각 클래스가 자신이 담당하는 초기화와 처리를 명확히 구분해야 한다.

- MRO의 끝에서 호출이 자연스럽게 종료되도록 구성해야 한다.
\end{dhksmry}

\subsubsection{다중 상속에서 \zsh{__init__()}은 어떻게 호출되는가?}

다음 코드를 살펴보자.

\begin{dhkcode}
class A:
    def __init__(self):
        print("A 초기화")


class B(A):
    def __init__(self):
        print("B 초기화 시작")
        super().__init__()
        print("B 초기화 종료")


class C(A):
    def __init__(self):
        print("C 초기화 시작")
        super().__init__()
        print("C 초기화 종료")


class D(B, C):
    def __init__(self):
        print("D 초기화 시작")
        super().__init__()
        print("D 초기화 종료")

d = D()
\end{dhkcode}

실행 결과:

\begin{dhkzsh}
D 초기화 시작
B 초기화 시작
C 초기화 시작
A 초기화
C 초기화 종료
B 초기화 종료
D 초기화 종료
\end{dhkzsh}

여기서 \zsh{A.__init__()}은 한 번만 호출된다.

하지만 그 이유를 정확하게 이해해야 한다.

Python이 모든 부모 클래스의 \zsh{__init__()}을 자동으로 한 번씩 호출하기 때문은 아니다.

각 클래스의 \zsh{__init__()}이 \zsh{super.__init__()}을 호출하고, 그 호출이 MRO를 따라 이어지도록 작성되었기 때문이다.

만약 중간 클래스에서 \zsh{super().__init__()}을 생략하면 그 뒤의 초기화 체인이 중단될 수 있다.

\begin{dhknote}
\begin{dhksmry}[centered]
object.__init__()은 무엇을 할까?
\end{dhksmry}
\end{dhknote}

MRO의 마지막에는 보통 \zsh{object}가 존재한다.

\zsh{object.__init__()}은 기본적인 초기화 메서드이며, 사용자 정의 속성을 설정하는 특별한 초기화 작업을 수행하지 않는다.

따라서 협력적 초기화 체인의 마지막 단계로 사용할 수 있다.

다만 \zsh{object.__init__()}은 임의의 사용자 정의 키워드 인수를 받아들이지 않으므로, 다중 상속에서 인수를 전달할 때는 각 클래스가 자신에게 필요한 인수를 처리하고 나머지를 다음 메서드로 전달하도록 설계해야 한다.

\subsubsection{\zsh{super()}와 일반적인 메서드 호출의 차이}

지금까지의 내용을 비교해 보자.

\begin{dhkcode}
class A:
    def method(self):
        print("A")


class B(A):
    def method(self):
        super().method()
\end{dhkcode}

여기서 다음 세 표현식은 서로 다른 의미를 가진다.

\begin{dhkcode}
self.method()
\end{dhkcode}

현재 인스턴스에서 일반적인 속성 탐색을 수행한다.

따라서 \zsh{B} 인스턴스에서 호출하면 보통 \zsh{B.method()}가 다시 선택된다.

\zsh{B.method()} 내부에서 \zsh{self.method()}를 호출한다면 재귀 호출이 발생할 수 있다.

\begin{dhkcode}
A.method(self)
\end{dhkcode}

\zsh{A} 클래스의 \zsh{method} 함수 객체를 명시적으로 선택해 호출한다.

이 경우 MRO에서 다음 클래스를 자동으로 선택하는 것은 아니다.

\begin{dhkcode}
super().method()
\end{dhkcode}

현재 메서드의 기준 클래스 다음부터 MRO를 따라 \zsh{method}를 탐색한다.

이 차이를 표로 정리하면 다음과 같다.

\begin{dhktable}
|표현식|	탐색 방식|
|:---|:---|
|self.method()|	인스턴스의 일반적인 속성 탐색|
|A.method(self)|	A 클래스의 함수를 명시적으로 선택|
|super().method()|	MRO에서 기준 클래스 다음부터 탐색|
\end{dhktable}

이 세 가지를 구분하면 상속 코드의 실행 흐름을 훨씬 정확하게 분석할 수 있다.

\subsubsection{PyTorch와의 연결 - \zsh{super().__init__()}의 의미}

이제 PyTorch 모델 코드를 다시 보자.

\begin{dhkcode}
import torch
import torch.nn as nn


class MyModel(nn.Module):
    def __init__(self):
        super().__init__()

        self.linear = nn.Linear(10, 5)

    def forward(self, x):
        return self.linear(x)
\end{dhkcode}

여기서 다음 표현식은

\begin{dhkcode}
super().__init__()
\end{dhkcode}

단순히 \zsh{nn.Module.__init__(self)}를 축약한 문법으로만 이해하면 부족하다.

\zsh{MyModel}의 MRO에서 \zsh{MyModel} 다음에 위치한 클래스의 \zsh{__init__()}을 탐색하는 호출이다.

일반적인 단일 상속 구조에서는 다음과 같다.

\begin{dhkcode}
print([cls.__name__ for cls in MyModel.__mro__])
\end{dhkcode}

실행 결과:

\begin{dhkzsh}
['MyModel', 'Module', 'object']
\end{dhkzsh}

따라서 이 경우 \zsh{super().__init__()}은 \zsh{nn.Module.__init__()}을 호출하게 된다.

그런데 만약 다른 클래스를 함께 상속하면 MRO가 달라질 수 있다.

예를 들어 다음과 같은 구조를 생각해 보자.

\begin{dhkcode}
class LoggingMixin:
    def __init__(self, *args, **kwargs):
        print("LoggingMixin 초기화")
        super().__init__(*args, **kwargs)


class MyModel(LoggingMixin, nn.Module):
    def __init__(self):
        super().__init__()

        self.linear = nn.Linear(10, 5)
\end{dhkcode}

이 클래스의 MRO는 다음과 같다.

\begin{dhksmry}[centered]
MyModel \ra LoggingMixin \ra nn.Module \ra object
\end{dhksmry}

따라서 \zsh{MyModel.__init__()}에서 \zsh{super().__init__()}을 호출하면 먼저 \zsh{LoggingMixin.__init__()}이 실행된다.

그다음 \zsh{LoggingMixin.__init__()} 안의 \zsh{super().__init__()}이 

\zsh{nn.Module.__init__()}을 호출하게 된다.

\begin{dhksmry}[centered]
MyModel.__init__()
\da
LoggingMixin.__init__()
\da
nn.Module.__init__()
\da
모델 내부 상태 초기화
\da
self.linear 등록
\end{dhksmry}

이처럼 \zsh{super()}는 상속 구조가 변경되더라도 MRO에 따라 다음 구현을 탐색할 수 있다.

다만 PyTorch에서 다중 상속을 설계할 때는 각 클래스의 초기화 방식이 협력적 상속을 지원하는지 확인해야 한다.

모든 클래스가 임의의 다중 상속 조합을 안전하게 지원하는 것은 아니다.

\subsubsection{핵심 정리}

이번 학습에서 반드시 이해해야 할 내용을 정리해 보자.

\begin{dhknote}
\begin{dhksmry}[centered]
MRO는 클래스 속성 탐색에 사용되는 선형 순서다.
\end{dhksmry}
\end{dhknote}

\begin{dhkcode}
Dog.__mro__
Dog.mro()
\end{dhkcode}

이를 통해 탐색 순서를 확인할 수 있다.

\begin{dhknote}
\begin{dhksmry}[centered]
super()는 부모 클래스 객체를 반환하지 않는다.
\end{dhksmry}
\end{dhknote}

MRO를 기준으로 속성을 탐색하는 \zsh{super} 프록시 객체를 반환한다.

\begin{dhknote}
\begin{dhksmry}[centered]
일반적인 인스턴스 메서드 안의 super()는 다음과 대응한다.
\end{dhksmry}
\end{dhknote}

\begin{dhkcode}
super(CurrentClass, self)
\end{dhkcode}

단, 이 대응은 정상적인 메서드 문맥에서 사용되는 인수 없는 \zsh{super()}를 설명하는 것이다.

\begin{dhknote}
\begin{dhksmry}[centered]
다중 상속에서는 직접 부모와 MRO의 다음 클래스가 다를 수 있다.
\end{dhksmry}
\end{dhknote}

\begin{dhkzsh}
D → B → C → A → object
\end{dhkzsh}

이 구조에서 \zsh{B} 내부의 \zsh{super()}는 \zsh{C}의 메서드를 선택할 수 있다.

\begin{dhknote}
\begin{dhksmry}[centered]
Python은 C3 선형화로 MRO를 결정한다.
\end{dhksmry}
\end{dhknote}

C3는 부모 클래스의 상대적 순서와 기존 상속 계층의 순서를 일관되게 유지하려고 한다.

\begin{dhknote}
\begin{dhksmry}[centered]
협력적 다중 상속은 super() 호출 체인을 활용한다.
\end{dhksmry}
\end{dhknote}

각 클래스가 MRO의 다음 구현으로 제어를 전달하도록 설계하는 방식이다.

\begin{dhknote}
\begin{dhksmry}[centered]
__init__()도 MRO의 영향을 받는다.
\end{dhksmry}
\end{dhknote}

부모 초기화가 자동으로 모두 실행되는 것은 아니며,

적절한 \zsh{super().__init__()} 호출이 필요하다.

\begin{dhknote}
\begin{dhksmry}[centered]
PyTorch의 super().__init__()도 같은 원리로 동작한다.
\end{dhksmry}
\end{dhknote}

일반적인 \zsh{MyModel(nn.Module)}에서는 \zsh{nn.Module.__init__()}을 호출하지만, 다중상속에서는 중간에 다른 클래스가 위치할 수 있다.

\subsubsection{2.9에서 반드시 이해해야 할 결론}

이번 학습의 핵심은 다음 구조다.

\[
\boxed{
\begin{gathered}
\text{Instance's MRO}\\
\downarrow\\
[C_0,C_1,\ldots,C_n]\\
\downarrow\\
\texttt{super}(C_i,\text{instance})\\
\downarrow\\
C_{i+1}\text{부터 속성 탐색}\\
\downarrow\\
\text{다음 구현 선택}
\end{gathered}
}
\]

가장 중요한 문장은 다음과 같다.

\begin{dhksmry}[centered]
super()는 직접적인 부모 클래스를 호출하는 문법이 아니라,

MRO에서 지정한 기준 클래스 다음부터 속성을 탐색하는 메커니즘이다.
\end{dhksmry}

따라서 같은 클래스에 정의된 \zsh{super()} 호출이라도 실제 인스턴스의 MRO가 달라지면 다음에 선택되는 메서드가 달라질 수 있다.

이 원리를 이해하면 Python의 다중 상속뿐 아니라 PyTorch의 \zsh{nn.Module}을 상속한 클래스에서 \zsh{super().__init__()}을 사용하는 이유도 더 정확하게 설명할 수 있다.

다음 2.10에서는 지금까지 학습한 객체 호출, 속성 탐색, 메서드 바인딩을 연결하여

\zsh{__call__, __len__, __getitem__, __iter__} 등의 특수 메서드가 Python 문법과 어떻게 연결되는지 살펴볼 것이다.
